\begin{tikzpicture}
\foreach \i in {0,...,7}{
  \node[note] at (1.3*\i,3.6) {\i};
  \node[draw=Blue,minimum width=1.1cm,minimum height=.7cm,fill=blue!3] at (1.3*\i,2.9) {0};}
\foreach \i in {1,4,6}{
  \node[draw=Blue,minimum width=1.1cm,minimum height=.7cm,fill=green!15] at (1.3*\i,2.9) {1};}
\node[box,text width=4cm] (k1) at (1.5,5.1) {加入 K1：位置 1、4};
\node[box,text width=4cm] (k2) at (7.7,5.1) {加入 K2：位置 4、6};
\draw[arr] (k1.south)--(1.3,3.3);
\draw[arr] (k1.south)--(5.2,3.3);
\draw[arr] (k2.south)--(5.2,3.3);
\draw[arr] (k2.south)--(7.8,3.3);
\node[warn,text width=4.4cm] (k3) at (1.6,.9) {查询 K3：位置 1、6 都为 1\\K3 未加入，得到假阳性};
\node[good,text width=4.4cm] (k4) at (7.6,.9) {查询 K4：位置 0、5 为 0\\对本插入集合判为未见};
\draw[arr,draw=Gold,dashed] (1.3,2.5)--(k3.north);
\draw[arr,draw=Gold,dashed] (7.8,2.5)--(k3.north);
\draw[arr,draw=green!35!black,dashed] (0,2.5)--(k4.north);
\draw[arr,draw=green!35!black,dashed] (6.5,2.5)--(k4.north);
\node[note,text width=12cm] at (4.6,-.4) {准确集合为 \{K1,K2\}；过滤器查询结果只决定下一步读取路径};
\end{tikzpicture}
